\section{Latar Belakang}
Surabaya merupakan salah satu kota metropolitan di Indonesia yang memiliki kepadatan aktivitas ekonomi, sosial, dan infrastruktur yang tinggi. \citep{saesarin2020skripsi}. Terletak di wilayah timur Pulau Jawa, Surabaya berada pada zona iklim tropis basah yang ditandai dengan adanya dua musim, yaitu musim penghujan dan musim kemarau. Kondisi iklim ini menjadikan curah hujan sebagai salah satu elemen penting dalam tata kelola perkotaan. Tingginya intensitas curah hujan dalam waktu singkat di Surabaya dapat menyebabkan genangan atau banjir yang berpotensi mengganggu aktivitas masyarakat dan merusak infrastruktur. \citep{febriana2019prediksi} 

Curah hujan sendiri merupakan salah satu parameter meteorologi yang paling krusial dalam berbagai sektor, seperti pertanian, pengelolaan sumber daya air, mitigasi bencana, serta perencanaan pembangunan wilayah. Curah hujan menunjukkan banyaknya air hujan yang turun dalam suatu wilayah selama periode waktu tertentu. Dalam konteks kota seperti Surabaya, pemahaman terhadap curah hujan sangat penting untuk menunjang kebijakan adaptif dan strategi mitigasi terhadap dampak bencana hidrometeorologi.\citep{azka2018gpm}

Dalam beberapa tahun terakhir terjadi ketidakpastian musim yang dipengaruhi oleh perubahan iklim global, termasuk pergeseran waktu dan intensitas musim penghujan maupun kemarau. Hal ini menjadikan upaya untuk memahami karakteristik curah hujan sebagai langkah penting dalam menghadapi tantangan perubahan iklim \citep{sofia2018merapi}. Salah satu karakteristik penting curah hujan adalah variabilitasnya, baik secara temporal (waktu) maupun spasial (lokasi geografis), yang mencerminkan fluktuasi jumlah serta distribusi hujan. 

Untuk mendukung perencanaan dan pengambilan kebijakan yang adaptif, diperlukan model prediksi curah hujan yang akurat. Pemodelan ini bertujuan untuk melihat dan menganalisis pola hujan yang akan terjadi di masa mendatang. Penelitian ini menggunakan data curah hujan berbasis satelit dari NASA untuk menganalisis dinamika curah hujan dan meramalkan pola curah hujan di wilayah Surabaya. Diharapkan hasil dari peramalan ini dapat digunakan sebagai dasar perencanaan kebijakan perkotaan yang tanggap terhadap risiko bencana dan perubahan iklim.

\textit{Forecasting} atau peramalan adalah proses membuat prediksi mengenai kejadian di masa depan dengan menggunakan data historis dan pengetahuan saat ini. Tujuannya adalah memberikan dasar yang kuat bagi pengambilan keputusan dalam berbagai bidang. Peramalan tidak selalu harus menghasilkan satu nilai pasti, dalam praktiknya dapat berupa nilai tunggal \textit{point forecast}, interval prediksi, distribusi probabilitas atau kemungkinan terjadinya sesuatu. Proses peramalan melibatkan pengamatan pola masa lampau seperti tren, musiman, atau siklus, yang kemudian dianalisis menggunakan berbagai metode statistik atau teknik pembelajaran mesin \citep{petropoulos2021forecasting}. Terdapat dua jenis metode peramalan yang terdiri dari kualitatif dan kuantitatif. Metode kualitatif mengandalkan pendapat ahli dan bersifat subjektif, biasanya digunakan saat data historis sulit diperoleh. Sebaliknya, metode kuantitatif menggunakan prinsip statistik, lebih objektif, dan mencakup berbagai pendekatan salah satunya yaitu \textit{time series} yang menganalisis pola data berurutan berdasarkan waktu untuk menghasilkan prediksi masa depan \citep{ruchjana2019spatiotemporal}.

Namun, metode peramalan konvensional \textit{time series} murni hanya mempertimbangkan aspek temporal tanpa mengikutsertakan faktor spasial, sehingga kurang optimal dalam menganalisis fenomena seperti curah hujan yang dipengaruhi oleh kondisi geografis wilayah. Itu dibuktikan dalam penelitian yang berjudul \textit{”Enhanced Spatio Temporal Modeling for Rainfall Forecasting: A High Resolution Grid Analysis”} yang menjelaskan bahwa model konvensional seperti ARIMA hanya menangkap aspek temporal dan linier. Tanpa mempertimbangkan korelasi spasial antar lokasi, model model tersebut sering kali gagal memberikan prediksi curah hujan yang akurat terutama ketika diterapkan pada wilayah dengan variasi geografis yang kompleks. Sebaliknya, pendekatan spatio temporal berbasis grid resolusi tinggi merupakan salah satu pendekatan yang relevan dalam konteks iklim dan meteorologi \citep{alam2024enhanced}.

Model spasio-temporal merupakan suatu pendekatan yang menggambarkan fenomena alam dengan mempertimbangkan aspek lokasi (spasial) dan waktu (temporal) \citep{utami2017gs3sls}. Dalam proses analisisnya, model ini memperhitungkan adanya ketergantungan antar lokasi pengamatan serta hubungan korelatif yang mungkin muncul dari keterlambatan waktu (lag). Data yang diamati berdasarkan waktu biasanya tidak bersifat independen, melainkan saling berkorelasi dan membentuk pola deret waktu. Model spasio-temporal pertama kali dikembangkan oleh Bilonick dan Nicholas pada tahun 1983, melalui penelitian curah hujan menggunakan tiga pendekatan. Pendekatan pertama berfokus pada aspek temporal dengan mengabaikan pengaruh spasial, yang kedua menekankan pengaruh spasial tanpa mempertimbangkan aspek waktu, dan yang ketiga menggabungkan keduanya dalam satu analisis \citep{islami2021spatiotemporal}.

Pendekatan spasio-temporal memungkinkan identifikasi kemiripan pola curah hujan antar wilayah dengan memperhatikan hubungan geografis antar lokasi, baik yang berdekatan maupun berjauhan. Salah satu alat penting dalam pendekatan ini adalah \textit{Spatial Weight Matrix}. \textit{Spatial Weight Matrix} merupakan matriks yang merepresentasikan kekuatan hubungan spasial antara suatu lokasi dengan lokasi lainnya \citep{mukhaiyar2024generalized}. Nilai dalam matriks ini menunjukkan seberapa besar pengaruh satu lokasi terhadap lokasi lain berdasarkan kedekatan geografis, jarak, atau karakteristik spasial tertentu.

Dalam pengembangannya, \textit{Spatial Weight Matrix} dapat dibentuk dalam dua dimensi (menggunakan koordinat horizontal) maupun tiga dimensi (menggabungkan perbedaan ketinggian). Pendekatan tiga dimensi ini menjadi penting ketika fenomena yang diteliti, seperti curah hujan atau muka air tanah, dipengaruhi oleh topografi atau elevasi wilayah. Matrix ini bisa dibentuk dengan berbagai metode seperti bobot berbasis inverse distance, bobot seragam, hingga yang dimodifikasi untuk memasukkan faktor ketinggian \citep{mukhaiyar2024generalized}.

Dengan menggunakan \textit{Spatial Weight Matrix} yang tepat, model spatio-temporal seperti STARMA \textit{(Space Time Autoregressive Moving Average)} mampu menangkap keterikatan spasial dan temporal secara simultan \citep{sharma2025integrating}. Model ini mengakomodasi efek pengaruh wilayah sekitar terhadap suatu lokasi, serta mempertimbangkan dinamika temporal yang mempengaruhi dari waktu ke waktu. STARMA menggunakan bobot dan mengintegrasikannya dengan temporal lag untuk menghasilkan prediksi yang lebih akurat terhadap fenomena spasio-temporal seperti curah hujan.  

Penggunaan pendekatan spasio-temporal seperti STARMA menjadi semakin relevan di tengah kondisi perubahan iklim yang menyebabkan curah hujan menjadi lebih ekstrem, tidak teratur, dan tidak menentu di berbagai wilayah Indonesia. Oleh karena itu, perancangan model prediksi curah hujan berbasis spasio-temporal tidak hanya memiliki kontribusi dalam pengembangan ilmu, tetapi juga sangat krusial untuk mendukung sistem peringatan dini, kebijakan mitigasi bencana, dan pengelolaan sumber daya air secara berkelanjutan.

Dalam mendukung kebutuhan tersebut, model spasio-temporal seperti STARMA telah dikembangkan dan digunakan untuk menangani fenomena yang melibatkan keterkaitan ruang dan waktu secara simultan. Konsep dasar dari model ini berakar pada penelitian awal oleh Cliff dan Ord pada tahun 1975 melalui pengembangan model STAR \textit{(Space-Time Autoregressive)} serta model lanjutannya, yaitu STARMA \textit{(Space-Time Autoregressive Moving Average)} \citep{laily2020hybrid}. Model STARMA menggabungkan pengaruh lag waktu sebelumnya baik dalam komponen autoregresif maupun \textit{moving average}, yang berlaku tidak hanya secara temporal tetapi juga secara spasial. Model ini digunakan untuk memodelkan data yang menunjukkan adanya autokorelasi spasial, yaitu keterkaitan antara nilai-nilai di lokasi yang berbeda dalam suatu sistem spasio-temporal \citep{rathod2018improved}.

Berbagai pendekatan telah digunakan dalam peramalan curah hujan, mulai dari model \textit{time series} konvensional seperti ARIMA yang menangani pola linier, hingga metode pembelajaran mesin seperti ANN, SVR, dan LSTM untuk menangani pola non-linier. Model \textit{hybrid machine learning} menjadi alternatif menjanjikan karena mampu meningkatkan akurasi prediksi melalui kombinasi teknik dekomposisi sinyal dan algoritma optimasi, meskipun keberhasilannya sangat bergantung pada arsitektur dan pemrosesan data yang tepat \citep{latif2022review}. 

Penelitian di Andhra Pradesh menggabungkan ARIMA, LSTM, XGBoost, dan Prophet dalam model ensemble menggunakan weighted averaging, menghasilkan prediksi yang lebih akurat dibandingkan model individu, dengan kontribusi besar dari ARIMA dan LSTM \citep{sravan2023leveraging}. Di Indonesia, penelitian menggunakan algoritma \textit{clustering} berbasis DTW menunjukkan bahwa klasifikasi iklim tidak hanya berdasarkan rata-rata, tetapi juga mempertimbangkan pola spasio-temporal \citep{suryana2023time}. Penelitian lain membandingkan beberapa algoritma \textit{machine learning} dan menemukan bahwa \textit{Random Forest Regressor} paling akurat dalam memprediksi curah hujan berdasarkan atribut suhu, kelembapan, dan tekanan udara \citep{jagtap2019rainfall}. 

Namun, model konvensional dan \textit{machine learning} sering kurang efektif untuk data spasio-temporal karena kompleksitas polanya. Untuk itu, model STARMA dirancang khusus menangani data tersebut dan terbukti efektif dalam berbagai studi klimatologis. Contohnya, model hibrida STARMA-ANN-SVM digunakan untuk memprediksi curah hujan tahunan di Bengal Barat \citep{saha2020hybrid}. Di Indonesia, model GSTAR juga telah diterapkan pada data curah hujan di Jawa Tengah dengan hasil yang menunjukkan pengaruh spasial yang signifikan. Sementara itu, pengembangan model STARMA II yang menggabungkan STARMA dan TDNN berhasil meningkatkan akurasi prediksi hasil panen padi dengan menggunakan matriks spasial berbasis KNN dan korelasi Pearson untuk menggambarkan \textit{Spatio Temporal Similarity} antar wilayah \citep{badu-apraku2021two}. 

Dengan melihat studi-studi terdahulu terbukti bahwa penggunaan model STARMA dan variannya mampu memberikan hasil prediksi yang lebih akurat pada data yang memiliki ketergantungan ruang dan waktu, khususnya dalam dominan meteorologi dan iklim.